\chapter{Hot Flushes Due to Menopause}
\index{Hot Flushes Due to Menopause}

\section*{Definition and Overview}
Hot flushes, also called hot flashes or vasomotor symptoms, are sudden sensations of intense heat that spread over the face, neck, and chest, often accompanied by flushing, sweating, and sometimes palpitations and anxiety. They are the most common symptom of the menopausal transition, caused by changing levels of oestrogen that destabilise the brain's temperature-regulating centre. When hot flushes occur at night they are called night sweats and can severely disrupt sleep. Most flushes settle within a few years of the final period, and effective treatments, including hormone replacement therapy and lifestyle measures, are available.

\section*{Epidemiology}
Up to 75\% of women experience hot flushes during the menopausal transition, and around 25\% have symptoms severe enough to affect daily functioning. Flushes typically begin in perimenopause and continue for an average of 4--7 years, though a minority of women have symptoms for a decade or longer. Night sweats are a common cause of menopausal sleep disturbance. Hot flushes are more common and more severe in women who are overweight, who smoke, and after surgical menopause. Most women can find effective relief.

\section*{Aetiology and Pathophysiology}
\begin{itemize}
    \item \textbf{Oestrogen fluctuation:} declining and fluctuating oestrogen levels affect the hypothalamus, which regulates body temperature
    \item \textbf{Narrowed thermoneutral zone:} the brain's temperature set-point range narrows, so small temperature changes trigger heat-loss responses
    \item \textbf{Pathophysiology:} blood vessels in the skin dilate, causing flushing and sweating, followed by a feeling of chill as heat is lost
    \item \textbf{Triggers:} hot drinks, alcohol, spicy food, caffeine, stress, and warm environments can provoke flushes
    \item \textbf{Surgical menopause:} sudden oestrogen withdrawal causes more severe symptoms
    \item \textbf{Hormonal treatment:} oestrogen replacement stabilises the thermoregulatory system
\end{itemize}

\section*{Clinical Features}
\begin{itemize}
    \item Sudden sensation of intense heat, usually starting in the chest or face and spreading
    \item Skin redness or flushing and visible sweating
    \item Palpitations, anxiety, or a feeling of breathlessness during the flush
    \item Night sweats with drenching perspiration and disturbed sleep
    \item Flushes last 1--5 minutes and occur several times daily to hourly
    \item Chills after the flush as the body cools
    \item Symptoms ease over years as menopause is completed
\end{itemize}

\section*{Differential Diagnosis}
\begin{itemize}
    \item Thyroid disease can cause heat intolerance and sweating
    \item Infections and fever cause flushing and sweats
    \item Anxiety and panic attacks can mimic flushes
    \item Medications and alcohol can cause flushing
    \item Carcinoid syndrome and other rare conditions cause flushing
    \item The menopausal context and hormone assessment clarify the diagnosis
\end{itemize}

\section*{When to Seek Medical Care}
Women with troublesome flushes or night sweats should discuss treatment options with their doctor. Seek review if flushes are accompanied by weight loss, fever, or other concerning symptoms.

\textbf{Contact your local emergency services for:}
\begin{itemize}
    \item Chest pain, palpitations with dizziness, or breathlessness
    \item Fainting or collapse
    \item High fever with sweating and confusion
    \item Severe headache or neurological symptoms
\end{itemize}

\section*{Diagnosis and Investigations}
\begin{itemize}
    \item \textbf{Clinical diagnosis:} typical flushes in a woman of menopausal age are usually diagnosed without tests
    \item \textbf{Hormone tests:} FSH measurement may help in younger women or unclear cases
    \item \textbf{Thyroid function tests:} to exclude thyroid disease
    \item \textbf{Symptom diary:} frequency and triggers are documented
    \item \textbf{Assessment of impact:} effect on sleep, work, and quality of life
    \item \textbf{Risk assessment:} for hormone therapy decisions
\end{itemize}

\section*{Management}
\begin{itemize}
    \item \textbf{Hormone replacement therapy:} the most effective treatment, relieving flushes in nearly all women
    \item \textbf{Vaginal oestrogen:} for genitourinary symptoms that often accompany flushes
    \item \textbf{Non-hormonal medicines:} certain antidepressants and other agents reduce flushes in women who cannot use HRT
    \item \textbf{Cognitive behavioural therapy:} has proven benefit for vasomotor symptoms and sleep
    \item \textbf{Lifestyle measures:} layered clothing, cool drinks, fans, and avoiding triggers
    \item \textbf{Weight management and exercise:} reduce symptom severity
    \item \textbf{Tapering treatment:} gradually reducing HRT when symptoms settle
\end{itemize}

\section*{Complications}
\begin{itemize}
    \item Sleep deprivation from night sweats, with daytime fatigue and irritability
    \item Reduced quality of life, work performance, and relationships
    \item Anxiety and low mood related to symptoms
    \item Avoidance of social situations because of flushing
    \item Long-term consequences of untreated menopause, including bone loss, are addressed separately
\end{itemize}

\section*{Prognosis and Outcomes}
Hot flushes are temporary for most women, resolving within 2--5 years of the final period, although a minority experience them for longer. HRT controls flushes effectively while it is used, and CBT and lifestyle measures help those who cannot or choose not to use hormones. With appropriate management, the impact of flushes on sleep and daily life is substantially reduced, and symptoms settle naturally over time.

\section*{Prevention and Self-Care}
\begin{itemize}
    \item Dress in layers and remove clothing as a flush begins
    \item Keep cool: fans, cool drinks, and a cool bedroom
    \item Identify and avoid personal triggers: hot drinks, alcohol, spicy food, and caffeine
    \item Manage stress with relaxation and mindfulness
    \item Maintain a healthy weight and exercise regularly
    \item Protect sleep with a consistent routine and cooling strategies
    \item Discuss hormone and non-hormone options with your doctor
    \item Use complementary approaches only with awareness of limited evidence
\end{itemize}

\section*{Sources and Further Reading}
The following reputable sources provide further information for healthcare students and patients seeking reliable, current advice about menopausal hot flushes.

\begin{itemize}
    \item Australasian Menopause Society. \textit{Hot flushes and night sweats information.} https://www.menopause.org.au/
    \item Jean Hailes for Women's Health. \textit{Managing hot flushes.} https://www.jeanhailes.org.au/
    \item Healthdirect. \textit{Menopause symptoms and hot flushes.} https://www.healthdirect.gov.au/menopause
    \item NHS. \textit{Menopause symptoms: hot flushes.} https://www.nhs.uk/conditions/menopause/
    \item Mayo Clinic. \textit{Hot flashes: causes and treatment.} https://www.mayoclinic.org/
    \item North American Menopause Society. \textit{Hot flash management.} https://www.menopause.org/
\end{itemize}
